\documentclass[12pt,a4paper]{article}
\usepackage{fontspec}
\usepackage{polyglossia}
\setmainlanguage{russian}
\setotherlanguage{english}
\setmainfont{DejaVu Serif}
\setsansfont{DejaVu Sans}
\setmonofont{DejaVu Sans Mono}
\usepackage{amsmath,amssymb}
\usepackage{geometry}
\usepackage{setspace}
\usepackage{hyperref}
\usepackage{booktabs}
\usepackage{indentfirst}
\usepackage{enumitem}

\geometry{left=3cm, right=1.5cm, top=2cm, bottom=2cm}
\onehalfspacing
\setlength{\parindent}{1.25cm}

\hypersetup{
  colorlinks=true,
  linkcolor=black,
  urlcolor=blue,
  citecolor=black
}

\begin{document}

%% ---------------------------------------------------------------
%%  Титульный лист
%% ---------------------------------------------------------------
\begin{titlepage}
\begin{center}
{\small Федеральное государственное автономное образовательное учреждение\\
высшего образования\\[4pt]
\textbf{НАЦИОНАЛЬНЫЙ ИССЛЕДОВАТЕЛЬСКИЙ\\
УНИВЕРСИТЕТ\\
«ВЫСШАЯ ШКОЛА ЭКОНОМИКИ»}\\[8pt]
Факультет компьютерных наук\\
Образовательная программа «Анализ больших данных»}

\vspace{3cm}

{\Large\textbf{МАГИСТЕРСКАЯ ДИССЕРТАЦИЯ}}

\vspace{1cm}

{\large\textbf{Прогнозирование валового регионального продукта (ВРП)\\
субъектов Российской Федерации с помощью ансамблевых\\
методов машинного обучения}}

\vspace{3cm}

\begin{flushright}
\begin{tabular}{l}
\textbf{Выполнил:}\\
Семенюченко Артём Геннадьевич\\
студент ФКН НИУ ВШЭ,\\
ОП «Анализ больших данных»\\[12pt]
\textbf{Научный руководитель:}\\
Паточенко Евгений Анатольевич\\
старший преподаватель ФКН НИУ ВШЭ,\\
академический руководитель\\
ОП «Анализ больших данных»
\end{tabular}
\end{flushright}

\vfill
Москва 2026
\end{center}
\end{titlepage}

%% ---------------------------------------------------------------
%%  Аннотация
%% ---------------------------------------------------------------
\newpage
\section*{Аннотация}

Настоящая работа посвящена разработке воспроизводимого подхода
к прогнозированию валового регионального продукта субъектов Российской
Федерации на основе панельных временных данных и ансамблевых методов
машинного обучения. Эмпирическая база~--- региональная панель за 2005--2024
годы; горизонт прогнозирования~--- один год вперёд (h\,=\,1). В роли целевой
переменной рассматривается логарифм реального ВРП. Признаковое пространство
объединяет лаги отклика, производственные и инвестиционные показатели,
индикаторы рынка труда, бюджетные параметры, признаки кризисных периодов
и характеристики региональной специализации.

Сопоставляются три базовые модели~--- наивный прогноз, AR(1) и Ridge
регрессия с фиксированными эффектами~--- и три ансамблевых алгоритма:
LightGBM, CatBoost и GPBoost. Последний выбран как метод, совмещающий
градиентный бустинг со случайными эффектами и позволяющий явно учитывать
панельную структуру данных. Качество оценивается на временно\'{й} валидации
по схеме expanding window по метрикам RMSE, MAE и out-of-sample~$R^2$.
Дополнительно проводятся ablation study по блокам признаков и два
robustness checks. Особое внимание уделено предотвращению leakage: все
preprocessing-операции обучаются исключительно на train-части каждого fold.

\smallskip\noindent
\textbf{Ключевые слова:} валовой региональный продукт, машинное обучение,
панельные временные ряды, ансамблевые модели, градиентный бустинг,
GPBoost, временна\'{я} валидация, SHAP.

%% ---------------------------------------------------------------
%%  Abstract
%% ---------------------------------------------------------------
\newpage
\section*{Abstract}

This thesis develops a reproducible pipeline for forecasting the Gross Regional
Product (GRP) of Russian Federation subjects using panel time-series data and
ensemble machine learning methods. The empirical base is a regional panel
covering 2005--2024; the forecast horizon is one year ahead ($h=1$). The target
variable is defined as the logarithm of real GRP. The feature space combines
lagged values of the target, industrial and investment indicators, labour market
statistics, fiscal variables, crisis dummies, and regional specialisation
characteristics.

The study compares three baseline models~--- naive forecast, AR(1), and Ridge
regression with fixed effects~--- and three ensemble algorithms: LightGBM,
CatBoost, and GPBoost. GPBoost is selected as a method that combines gradient
boosting with Gaussian random effects, explicitly addressing panel heterogeneity
without expanding the feature space. Model quality is assessed via expanding-window
time-series cross-validation using RMSE, MAE, and out-of-sample~$R^2$. The
experimental design additionally includes an ablation study across feature
groups and two robustness checks. Particular attention is paid to leakage
prevention: all preprocessing operations are fitted exclusively on the training
portion of each fold.

\smallskip\noindent
\textbf{Keywords:} gross regional product, machine learning, panel time series,
ensemble methods, gradient boosting, GPBoost, time-series validation, SHAP.

%% ---------------------------------------------------------------
%%  Содержание
%% ---------------------------------------------------------------
\newpage
\tableofcontents

%% ---------------------------------------------------------------
%%  1. ВВЕДЕНИЕ
%% ---------------------------------------------------------------
\newpage
\section{Введение}

\subsection{Актуальность исследования}

Валовой региональный продукт занимает центральное место среди показателей
экономического развития субъектов Российской Федерации. По экономическому
содержанию ВРП близок к ВВП, но рассчитывается для отдельной территории,
что позволяет перейти от агрегированного описания к анализу пространственного
распределения выпуска и неоднородной реакции регионов на экономические шоки.

Российская экономика характеризуется выраженной территориальной
неоднородностью, обусловленной историей экономического становления
страны~--- централизованностью вокруг крупных городов и большим количеством
моногородов~--- а также особенностями ландшафта и климата (регионы-курорты,
добывающие регионы). В результате в одной стране сосуществуют сырьевые,
индустриальные, аграрные и бюджетно-зависимые регионы, каждый из которых
обладает собственным механизмом формирования валового продукта.
Прикладную значимость исследования усиливает лаг публикации официальной
статистики: по данным Росстата, объём ВРП по сумме субъектов за 2024 год
в текущих ценах составил 186\,545,1 млрд руб., индекс физического
объёма~--- 104,4\%~\cite{rosstat2024}. Эти оценки появляются после завершения
отчётного периода, тогда как региональный мониторинг и бюджетное планирование
требуют более оперативной картины динамики выпуска.

Сама задача относится к классу регрессии на панельных временных данных.
Наблюдением выступает пара «регион~--- год», целевой переменной~--- будущее
значение ВРП, признаками~--- лаговая информация и социально-экономические
показатели региона. Такая структура отличается от классического одномерного
временного ряда: модель должна одновременно учитывать временну\'{ю} зависимость,
межрегиональную неоднородность, различие масштабов субъектов и возможные
структурные разрывы. Дополнительная сложность состоит в риске утечки
информации из будущего: применение стандартных процедур кросс-валидации
без учёта временного порядка систематически завышает оценки качества.
Cerqua, Letta и Pinto~\cite{cerqua2025} показывают, что leakage в панельных
данных имеет несколько источников, не сводимых к простому «подглядыванию»
в тестовую выборку.

Современная литература по macro-ML подчёркивает, что применение машинного
обучения оправдано лишь при строгой экспериментальной процедуре~\cite{chi2025}.
Соответственно, ансамблевые методы рассматриваются в данной работе не как
заведомо более точные, а как гипотеза, подлежащая проверке относительно
сильных базовых моделей.

\subsection{Объект и предмет исследования}

\textit{Объектом} исследования выступают панельные временные ряды
социально-экономических показателей субъектов Российской Федерации.
\textit{Предметом}~--- методы прогнозирования ВРП на этих данных,
охватывающие как эконометрические базовые модели, так и ансамблевые
алгоритмы машинного обучения.

\subsection{Цель и задачи исследования}

\textit{Цель работы}~--- разработка и оценка воспроизводимого пайплайна
прогнозирования ВРП субъектов Российской Федерации на панельных временных
данных с использованием эконометрических базовых моделей и ансамблевых
методов машинного обучения.

Для достижения цели решаются следующие задачи:
\begin{enumerate}
  \item Провести обзор литературы по макроэкономическому и региональному
        прогнозированию, включая современные ML-подходы.
  \item Сформировать панельную базу данных за 2005--2024 годы, проверить
        стационарность и определить целевую переменную.
  \item Реализовать предобработку данных и leakage-free процедуру
        формирования признаков.
  \item Построить базовые модели (наивный прогноз, AR(1), Ridge с FE)
        и ансамблевые алгоритмы (LightGBM, CatBoost, GPBoost).
  \item Реализовать временну\'{ю} валидацию по схеме expanding window
        и подбор гиперпараметров строго внутри train-части каждого fold.
  \item Сравнить модели по RMSE, MAE и out-of-sample~$R^2$.
  \item Провести ablation study по блокам признаков и robustness checks.
  \item Построить SHAP-интерпретацию лучшей модели.
\end{enumerate}

\subsection{Данные и инструментарий}

Информационная база формируется на основе данных Росстата, ЕМИСС,
Банка России, Минэкономразвития и Федерального казначейства за 2005--2024
годы. В признаковое пространство включены лаги ВРП, индекс промышленного
производства, показатели добычи и обработки, сельское хозяйство,
строительство, транспорт, торговля, инвестиции в основной капитал, занятость,
безработица, заработная плата, доходы населения, доходы и расходы
региональных бюджетов, трансферты, федеральный округ и dummy-переменные
кризисных периодов.

Используемые библиотеки: \texttt{pandas~2.2}, \texttt{numpy~1.26}~--- работа
с данными; \texttt{scikit-learn~1.4}~--- линейные модели, валидация,
масштабирование; \texttt{lightgbm~4.3}, \texttt{catboost~1.2}~---
бустинговые алгоритмы; \texttt{gpboost~1.5}~--- модель со смешанными
эффектами; \texttt{linearmodels~6.0}~--- панельные регрессии с
фиксированными эффектами; \texttt{optuna~3.6}~--- подбор гиперпараметров;
\texttt{shap~0.45}~--- интерпретация; \texttt{matplotlib~3.8},
\texttt{seaborn~0.13}~--- визуализация. Реализация выполнена на Python~3.11.
Воспроизводимость обеспечивается фиксацией \texttt{random\_state=42} во всех
стохастических процедурах и сохранением конфигураций в YAML-файлах.

\subsection{Практическая значимость}

Результаты могут использоваться для мониторинга региональной экономической
динамики, предварительной оценки ВРП до публикации официальной статистики
и построения сценарных прогнозов. SHAP-разложения позволяют сопоставить
выявленные факторы с экономической логикой регионов, что делает модель
пригодной для содержательного аналитического использования.

\subsection{Научная новизна}

В существующих российских исследованиях акцент делается либо на
детерминантах ВРП~\cite{badykova2026}, либо на сценарном прогнозе для
отдельного региона~\cite{malkina2025}. Новизна настоящей работы~--- в
формировании единого экспериментального фреймворка: прогнозирование ВРП
как регрессия на панельных временных данных с временно\'{й} валидацией,
явным контролем leakage, ablation study, robustness checks и
SHAP-интерпретацией на единой панели субъектов. Дополнительно в сравнение
включён GPBoost~--- алгоритм, совмещающий градиентный бустинг со случайными
эффектами и ранее не применявшийся к задаче прогнозирования ВРП российских
регионов.

%% ---------------------------------------------------------------
%%  2. ОБЗОР ЛИТЕРАТУРЫ
%% ---------------------------------------------------------------
\section{Обзор литературы}

Задача прогнозирования ВРП пересекается с несколькими исследовательскими
направлениями: макроэкономическое прогнозирование, машинное обучение на
панельных данных и интерпретируемые модели. Ниже детально разбираются
четыре работы, наиболее значимых для методологии настоящего исследования,
а затем кратко резюмируется сопутствующая литература.

\subsection{Chi et al.\ (2025). Macroeconomic Forecasting and
Machine Learning~\cite{chi2025}}

\textit{Постановка.} В рамках данной работы акцент смещается с точечного
прогнозирования среднего на условное распределение макроэкономических
результатов с помощью машинного обучения. Рассматривается задача
прогнозирования уровня безработицы с длительным горизонтом. Методологически
статья объединяет три элемента современной прогнозной практики: работу с
высокоразмерными данными, регуляризацию для контроля сложности модели и
строгую out-of-sample валидацию.

\textit{Метод и ключевые результаты.} В статье используются глубокие
нейронные сети для прогнозирования квантилей распределения. В качестве
функции потерь применяется pinball loss (квантильная регрессия). Авторы
применяют expanding-window схему: модель обучается только на прошлых
наблюдениях, гиперпараметры подбираются на валидационном периоде, качество
проверяется на тестовом периоде. Главный результат: регуляризация играет
решающую роль. Без shrinkage высокоразмерная модель извлекает шум, а не
устойчивый макроэкономический сигнал. Нелинейные преобразования дают лишь
ограниченный прирост точности~--- основную роль играет дисциплинированная
процедура отбора сложности и корректная out-of-sample валидация.

\textit{Значимость для настоящей работы.} Статья задаёт методологический
стандарт для применения ML в макроэкономике: сложные модели сравниваются
с базовыми только во вневыборочной процедуре. Вывод об ограниченной роли
нелинейности предостерегает от автоматического ожидания превосходства
ансамблевых методов над линейными.

\subsection{Badykova (2026). Детерминанты ВРП в России:
подход машинного обучения~\cite{badykova2026}}

\textit{Постановка.} Автор анализирует детерминанты ВРП на душу населения
по 85 регионам России за 2013--2023 годы. Задача поставлена как объяснение
(feature importance), а не как прогнозирование.

\textit{Метод и результаты.} Основным алгоритмом выступает LightGBM.
Интерпретация построена на SHAP: глобальный анализ показывает важность
признаков по всей выборке, локальный~--- объясняет отклонения конкретных
регионов. Наиболее значимыми предикторами оказались инвестиции в основной
капитал и занятость в обрабатывающей промышленности.

\textit{Отличие от настоящей работы.} Badykova решает задачу объяснения,
а не прогнозирования. Настоящая работа переориентирована на вневыборочный
прогноз с временно\'{й} валидацией, tuning только внутри train, ablation
study и явным контролем leakage. Кроме того, в сравнение добавлен GPBoost.

\subsection{Lashina \& Grishunin (2023). Сравнение моделей
прогнозирования ВВП России~\cite{lashina2023}}

\textit{Постановка.} Авторы сопоставляют прогностическую точность четырёх
классов моделей применительно к ВВП России в условиях кризисных периодов:
ARIMA, VAR, SVR и CatBoost. Исследование охватывает кризисы 2014--2015
и 2020 годов.

\textit{Ключевые результаты.} CatBoost демонстрирует наилучшее качество
в кризисные периоды, превосходя ARIMA и VAR по RMSE и MAPE. В стабильные
периоды разница сокращается. Работа является единственной в рассмотренной
литературе, сравнивающей CatBoost с классическими эконометрическими
моделями на данных ВВП России и охватывающей кризисные периоды.

\textit{Значимость.} Результаты обосновывают включение CatBoost в сравнение
и мотивируют отдельный анализ качества моделей в кризисные годы (robustness
check).

\subsection{Sigrist (2022). Gaussian Process Boosting~\cite{sigrist2022}}

\textit{Постановка.} Работа предлагает алгоритм GPBoost, объединяющий
градиентный бустинг с гауссовыми процессами и моделями смешанных эффектов.
Мотивация: стандартный бустинг моделирует только условное среднее
$\mathbb{E}[y \mid X]$, тогда как в панельных данных остатки систематически
коррелированы внутри групп (регионов), что снижает эффективность оценки.

\textit{Модель.} Совместная спецификация:
\begin{equation}
  y_{i,t} = F(X_{i,t}) + b_i + \varepsilon_{i,t},
  \quad b_i \sim \mathcal{N}(0,\sigma_b^2),
  \quad \varepsilon_{i,t} \sim \mathcal{N}(0,\sigma^2),
\end{equation}
где $F(\cdot) = \sum_m \eta T_m(\cdot)$~--- аддитивная функция деревьев,
$b_i$~--- случайный эффект региона~$i$.

\textit{Значимость для настоящей работы.} GPBoost включён как алгоритм,
специально предназначенный для панельной структуры. Это позволяет проверить,
даёт ли явное моделирование региональной гетерогенности прирост качества
сверх стандартного бустинга.

\subsection{Сопутствующая литература}

\textit{Наукастинг.} Цукарев и соавторы~\cite{tsukarev2025} сравнивают
методы bridge equations и MIDAS с ML-алгоритмами на данных Армении
и Беларуси, показывая преимущество бустинга на агрегированных годовых данных.
Макеева и Станкевич~\cite{makeeva2022} применяют MIDAS и MFBVAR для
оперативной оценки ВВП России.

\textit{Региональные исследования.} Малкина, Сочков и Капустина~\cite{malkina2025}
разрабатывают нейросетевые сценарные прогнозы ВРП Нижегородской области.
Банк России~\cite{semiturkin2022} показывает, что целесообразность ML
следует проверять отдельно для каждой группы регионов.

\textit{Интерпретируемость и ML-дизайн.} Lundberg и Lee~\cite{lundberg2017}
разработали SHAP-разложение; в настоящей работе SHAP применяется для
проверки экономической согласованности модели, а не как инструмент
причинной идентификации. Смирнов~\cite{smirnov2025} предупреждает, что
прирост ML относительно классических моделей на традиционных данных
нередко умеренный~--- этот тезис усиливает необходимость строгого сравнения
с наивным прогнозом и AR(1).

%% ---------------------------------------------------------------
%%  3. МЕТОДОЛОГИЯ
%% ---------------------------------------------------------------
\section{Методология исследования}

\subsection{Формальная постановка задачи}

Пусть $i \in \{1,\ldots,N\}$~--- субъект Федерации, $t \in \{1,\ldots,T\}$~---
год, $h=1$~--- горизонт прогноза. Задача состоит в нахождении функции
$f_\theta$ такой, что
\begin{equation}
  \hat{y}_{i,t+1} = f_\theta(X^{\text{available}}_{i,t}),
\end{equation}
где $y_{i,t} = \log(\mathrm{GRP}^{\text{real}}_{i,t})$~--- логарифм реального
ВРП, $X^{\text{available}}_{i,t}$~--- вектор признаков, содержащий только
информацию, доступную к моменту~$t$. Множество данных:
\begin{equation}
  \mathcal{D} = \{(X_{i,t}, y_{i,t+1}) : i = 1,\ldots,N;\; t = p,\ldots,T-1\},
\end{equation}
где $p$~--- максимальный лаг признаков.

\subsection{Проверка стационарности}

Перед построением моделей проводится проверка стационарности. Для каждого
региона применяется расширенный тест Дики~--~Фуллера (ADF) на логарифм ВРП
и его первую разность. На уровне панели используется тест Im~--~Pesaran~--~Shin
(IPS), допускающий неодинаковые авторегрессионные коэффициенты по регионам.
Гипотеза о нестационарности уровня ВРП ожидаемо не отвергается; первые
разности предположительно стационарны. Эти результаты обосновывают
включение лагов ВРП как признаков.

\subsection{Признаковое пространство}

Признаковая система строится как экономически мотивированная структура:
\begin{equation}
  X_{i,t} = \bigl[X^{\text{lag}}_{i,t},\; X^{\text{prod}}_{i,t},\;
  X^{\text{inv}}_{i,t},\; X^{\text{lab}}_{i,t},\; X^{\text{bud}}_{i,t},\;
  X^{\text{cat}}_{i,t},\; X^{\text{shock}}_t\bigr].
\end{equation}

\textit{Инерционный блок} объединяет лаги ВРП $y_{i,t-1}$, $y_{i,t-2}$
и скользящие средние. \textit{Производственный}~--- показатели добычи,
обработки, промышленности, сельского хозяйства, строительства, транспорта
и торговли. \textit{Инвестиционный}~--- инвестиции в основной капитал,
инвестиции на душу населения, отношение инвестиций к ВРП.
\textit{Социально-трудовой}~--- занятость, безработица, заработная плата,
доходы населения. \textit{Бюджетный}~--- доходы и расходы бюджета,
трансферты, бюджетные инвестиции. \textit{Категориальный блок}~---
федеральный округ и индикаторы специализации (сырьевая, индустриальная,
аграрная, сервисная, бюджетно-зависимая). \textit{Блок шоков}~--- dummy
кризисных периодов 2008--2009, 2014--2015, 2020 и 2022--2023 годов.

\subsection{Предотвращение leakage и предобработка}

Применительно к панельным данным leakage имеет несколько
источников~\cite{cerqua2025}. Базовое правило:
\begin{equation}
  X^{\text{available}}_{i,t} = \{x^{(j)}_{i,s} : s \le t - \ell_j,\;
  j = 1,\ldots,p\},
\end{equation}
где $\ell_j \ge 1$~--- лаг доступности признака. Все операции нормализации
и импутации обучаются исключительно на $\mathrm{Train}_k$ и применяются
к $\mathrm{Test}_k$. Денежные показатели дефлируются индексами
потребительских цен. Для масштабирования применяется \texttt{RobustScaler}:
\begin{equation}
  x_{\text{scaled}} = \frac{x - Q_2}{Q_3 - Q_1},
\end{equation}
где $Q_1, Q_2, Q_3$~--- квартили обучающей выборки. Выбор \texttt{RobustScaler}
перед \texttt{StandardScaler} обусловлен структурными выбросами
(ХМАО, ЯНАО, Москва), не являющимися ошибками измерения, но способными
смещать оценку дисперсии.

\subsection{Базовые модели}

В работе используются три базовые модели, представляющие нарастающую
сложность спецификаций.

\textit{Наивный прогноз} воспроизводит последнее известное значение:
$\hat{y}_{i,t+1} = y_{i,t}$.

\textit{AR(1) по регионам} оценивается отдельно для каждого субъекта:
$y_{i,t+1} = \alpha_i + \rho_i y_{i,t} + \varepsilon_{i,t+1}$.

\textit{Ridge с фиксированными эффектами} объединяет индивидуальные
константы региона и временны\'{х} периодов с L2-регуляризацией полного
набора признаков:
\begin{equation}
  \hat{\beta}^{\text{ridge}} = \arg\min_{\beta}
  \Bigl[\sum_{i,t}(y_{i,t+1} - \alpha_i - \delta_t - X^\top_{i,t}\beta)^2
  + \lambda\|\beta\|^2_2\Bigr].
\end{equation}

\subsection{Ансамблевые модели}

В работе сравниваются три ансамблевых алгоритма.

\textit{LightGBM}~\cite{ke2017} реализует градиентный бустинг на
деревьях с leaf-wise стратегией роста, эксклюзивным разбиением признаков
(EFB) и гистограммным методом. Алгоритм выбран как основной представитель
класса gradient boosting ввиду высокой вычислительной эффективности.

\textit{CatBoost}~\cite{dorogush2018} применяет ordered boosting для
устранения target leakage при обработке категориальных признаков и
использует симметричные деревья. По результатам Lashina \& Grishunin~\cite{lashina2023}
демонстрирует наилучшее качество на российских данных в кризисные периоды.

\textit{GPBoost}~\cite{sigrist2022} совмещает бустинг со случайными
эффектами регионов:
\begin{equation}
  y_{i,t} = F(X_{i,t}) + b_i + \varepsilon_{i,t},
  \quad b_i \sim \mathcal{N}(0,\sigma_b^2),
  \quad \varepsilon_{i,t} \sim \mathcal{N}(0,\sigma^2).
\end{equation}
В отличие от включения dummy-переменных регионов, GPBoost оценивает
случайные эффекты ковариационно, не увеличивая размерность признакового
пространства и позволяя строить прогнозы для регионов с короткой историей
наблюдений.

\subsection{Временна\'{я} валидация и подбор гиперпараметров}

Основная схема~--- expanding window:
\begin{equation}
  \mathrm{Train}_k = \{(i,t) : t \le \tau_k\},\quad
  \mathrm{Test}_k = \{(i,t) : t = \tau_k + 1\}.
\end{equation}
Начальное окно обучения~--- не менее десяти лет (2005--2014),
$\tau_k$ последовательно увеличивается. Итоговый размер~--- восемь--десять
fold'ов.

Подбор гиперпараметров выполняется через дополнительное внутреннее
разбиение $\mathrm{Train}_k = \mathrm{Train}^{\text{inner}}_k \cup
\mathrm{Val}^{\text{inner}}_k$ с условием
$\max(\mathrm{Train}^{\text{inner}}_k) < \min(\mathrm{Val}^{\text{inner}}_k)$.
Параметры выбираются через Optuna (TPE sampler), после чего модель
переобучается на полном $\mathrm{Train}_k$.

\subsection{Метрики качества}

Основные метрики для сравнения моделей:
\begin{align}
  \mathrm{RMSE} &= \sqrt{\frac{1}{n}\sum_j(y_j - \hat{y}_j)^2}, &
  \mathrm{MAE} &= \frac{1}{n}\sum_j|y_j - \hat{y}_j|, \\[4pt]
  R^2_{\mathrm{OOS}} &= 1 - \frac{\sum_j(y_j-\hat{y}_j)^2}
  {\sum_j(y_j-\hat{y}^{\text{bench}}_j)^2},
\end{align}
где в знаменателе $R^2_{\mathrm{OOS}}$ в качестве benchmark используется
наивный прогноз. Положительный $R^2_{\mathrm{OOS}}$ означает, что модель
превосходит наивный прогноз за пределами обучающей выборки.

%% ---------------------------------------------------------------
%%  4. ПЛАН ЭКСПЕРИМЕНТОВ
%% ---------------------------------------------------------------
\section{План экспериментов и интерпретации}

\subsection{Ablation study}

Ablation study оценивает вклад отдельных блоков признаков через четыре
сценария:
\begin{itemize}
  \item \textbf{A0}~--- только лаги ВРП (минимальная спецификация);
  \item \textbf{A1}~--- лаги + производственный и инвестиционный блоки;
  \item \textbf{A2}~--- лаги + социально-трудовой и бюджетный блоки;
  \item \textbf{A3}~--- полный набор признаков, включая региональную
                        типологию и кризисные dummy.
\end{itemize}
Сравнение A0~--~A3 позволяет установить, дают ли производственные,
инвестиционные и бюджетные признаки прирост точности сверх лагов ВРП.

\subsection{Robustness checks}

Проводятся два robustness check.

\textit{Первый~--- исключение структурных выбросов.} Из обучающей и
тестовой выборок исключаются Москва, ХМАО, ЯНАО и Сахалинская область.
Цель~--- проверить устойчивость ранжирования моделей при удалении
регионов, доминирующих по масштабу ВРП.

\textit{Второй~--- кризисные против докризисных периодов.} Оценка
качества раздельно проводится на fold'ах, соответствующих кризисным годам
(2008--2009, 2014--2015, 2020, 2022--2023), и на fold'ах вне этих периодов.
Это позволяет верифицировать тезис Lashina \& Grishunin~\cite{lashina2023}
о преимуществе бустинга в периоды структурных разрывов.

\subsection{SHAP-интерпретация}

Для лучшей модели строится SHAP-разложение:
\begin{equation}
  \hat{y}_{i,t} = \phi_0 + \sum_{j=1}^{p} \phi_{j,i,t}.
\end{equation}
Глобальный анализ: $I_j = \frac{1}{n}\sum_{i,t}|\phi_{j,i,t}|$.
Групповой: $I_{j,g} = \frac{1}{|g|}\sum_{(i,t)\in g}|\phi_{j,i,t}|$.
Локальный~--- для отдельных субъектов. SHAP применяется как инструмент
объяснения модели, а не причинной идентификации.

\subsection{Анализ ошибок}

По регионам: $\mathrm{MAE}_i = \frac{1}{T}\sum_t|y_{i,t}-\hat{y}_{i,t}|$.
По годам: $\mathrm{MAE}_t = \frac{1}{N}\sum_i|y_{i,t}-\hat{y}_{i,t}|$.
Анализ позволяет выявить систематические ошибки для определённых типов
регионов и временны\'{х} периодов.

\subsection{Ожидаемые результаты и ограничения}

Ожидается, что лаги ВРП окажутся наиболее сильными предикторами; ключевой
вопрос~--- дают ли производственные, инвестиционные и бюджетные признаки
прирост сверх лагов. Не исключён сценарий, при котором Ridge с FE оказывается
сопоставим с бустинговыми алгоритмами~--- это означало бы, что относительно
короткая панель ограничивает выгоду от нелинейных алгоритмов.

Основные ограничения: (1)~около двадцати годовых наблюдений на регион;
(2)~ретроспективная база не воспроизводит real-time environment;
(3)~структурные разрывы затрудняют перенос модели между режимами;
(4)~типология регионов условна; (5)~SHAP не доказывает причинности.

%% ---------------------------------------------------------------
%%  5. ДАЛЬНЕЙШИЙ ПЛАН РЕАЛИЗАЦИИ
%% ---------------------------------------------------------------
\section{Дальнейший план реализации}

Настоящий раздел описывает конкретные шаги, необходимые для завершения
исследования от текущего состояния (методологический фреймворк описан,
данные частично собраны) до финальной защиты.

\subsection*{Этап 1. Сборка и верификация панели данных (2--3 недели)}

Необходимо завершить формирование панельной базы данных за 2005--2024 годы
по всем субъектам Российской Федерации. Конкретные шаги:
\begin{enumerate}
  \item Выгрузить данные из ЕМИСС по всем включённым показателям
        (промышленность, инвестиции, занятость, бюджет).
  \item Пересчитать номинальный ВРП в реальный с использованием
        региональных дефляторов Росстата. При их отсутствии~---
        использовать общероссийский ИПЦ как консервативную замену.
  \item Проверить наполненность панели: доля пропусков по каждому признаку
        не должна превышать 20\%. Признаки с бо\'{о}льшей долей пропусков
        исключить или заменить агрегированным прокси.
  \item Провести тест ADF и IPS. Зафиксировать результаты в отдельной
        таблице (включить в приложение к диссертации).
\end{enumerate}

\subsection*{Этап 2. Реализация пайплайна и базовых моделей (1--2 недели)}

\begin{enumerate}
  \item Реализовать модуль leakage-free предобработки с интерфейсом
        \texttt{sklearn.Pipeline}.
  \item Реализовать expanding-window валидацию как итератор fold'ов.
  \item Обучить и оценить три базовые модели: наивный прогноз, AR(1),
        Ridge с FE. Убедиться, что RMSE наивного прогноза воспроизводит
        ожидаемые значения, прежде чем переходить к ансамблям.
  \item Зафиксировать baseline-результаты в сводной таблице метрик.
\end{enumerate}

\subsection*{Этап 3. Обучение ансамблевых моделей (1--2 недели)}

\begin{enumerate}
  \item Обучить LightGBM с подбором гиперпараметров через Optuna
        (не более 50 trials на fold).
  \item Обучить CatBoost аналогично, с учётом категориальных признаков.
  \item Обучить GPBoost: начать с базовой спецификации случайных эффектов
        по регионам, при необходимости добавить временны\'{е} эффекты.
  \item Занести все результаты в единую таблицу: RMSE, MAE,
        $R^2_{\mathrm{OOS}}$ по каждому fold и в среднем.
\end{enumerate}

\subsection*{Этап 4. Ablation study и robustness checks (1 неделя)}

\begin{enumerate}
  \item Провести ablation study (сценарии A0~--~A3) для лучшей из
        ансамблевых моделей.
  \item Провести два robustness check: без структурных выбросов
        и отдельно для кризисных периодов.
\end{enumerate}

\subsection*{Этап 5. SHAP-интерпретация (3--5 дней)}

\begin{enumerate}
  \item Построить global SHAP importance для лучшей модели.
  \item Построить групповой SHAP по типам регионов (сырьевые,
        индустриальные, аграрные, бюджетно-зависимые).
  \item Провести локальный анализ для 3--5 репрезентативных регионов.
  \item Верифицировать экономическую интерпретируемость: убедиться,
        что знаки SHAP-ценностей согласуются с экономической теорией.
\end{enumerate}

\subsection*{Этап 6. Написание главы «Результаты» и финальная редакция (1--2 недели)}

\begin{enumerate}
  \item Написать главу с описанием полученных результатов, включая
        сводную таблицу метрик, графики SHAP и анализ ошибок.
  \item Сформулировать основной вывод: при каких условиях ансамблевые
        модели превосходят базовые, и какова практическая величина
        прироста точности.
  \item Привести раздел «Заключение» в соответствие с фактическими
        результатами.
  \item Подготовить презентацию к предзащите (8--10 слайдов):
        постановка задачи, данные, методология, результаты, выводы.
\end{enumerate}

\subsection*{Итоговый временно\'{й} план}

\begin{center}
\begin{tabular}{lll}
\toprule
Этап & Содержание & Срок \\
\midrule
1 & Сборка и верификация панели & 2--3 нед. \\
2 & Пайплайн и базовые модели  & 1--2 нед. \\
3 & Ансамблевые модели          & 1--2 нед. \\
4 & Ablation study и robustness & 1 нед.    \\
5 & SHAP-интерпретация           & 3--5 дн.  \\
6 & Результаты и финальная правка & 1--2 нед. \\
\midrule
  & \textbf{Итого}              & \textbf{7--11 нед.} \\
\bottomrule
\end{tabular}
\end{center}

%% ---------------------------------------------------------------
%%  6. ЗАКЛЮЧЕНИЕ
%% ---------------------------------------------------------------
\section{Заключение}

В работе задача прогнозирования ВРП субъектов Российской Федерации
сформулирована как регрессия на панельных временных данных с горизонтом
$h=1$. Экономическая часть определяет предметную область, выбор признаков
и критерии интерпретации; техническая часть задаёт воспроизводимый
pipeline, процедуры предотвращения leakage, временну\'{ю} валидацию и
сопоставление шести моделей: наивного прогноза, AR(1), Ridge с FE,
LightGBM, CatBoost и GPBoost.

Ключевая методологическая позиция: ансамблевые методы не являются
самоцелью. Их применение должно быть оправдано устойчивым улучшением
прогноза относительно базовых моделей. Отличительная черта исследования~---
включение GPBoost, совмещающего бустинг со случайными эффектами регионов,
что позволяет напрямую адресовать панельную гетерогенность без
расширения признакового пространства.

Сокращение модельного меню по сравнению с первоначальным планом~--- с
двенадцати до шести алгоритмов, ограничение горизонта прогнозирования
одним годом и снижение числа robustness checks~--- обосновано
размером эмпирической базы: при около двадцати годовых наблюдениях
на регион широкое модельное сравнение сопряжено с риском
переподгонки выводов к конкретным разбивкам данных.

Следующий содержательный шаг~--- реализация описанного в разделе~5
плана экспериментов и формулировка основного вывода о сравнительных
преимуществах ансамблевых методов применительно к панели субъектов
Российской Федерации.

%% ---------------------------------------------------------------
%%  Список литературы
%% ---------------------------------------------------------------
\newpage
\begin{thebibliography}{20}

\bibitem{rosstat2024}
Федеральная служба государственной статистики. Валовой региональный
продукт за 2024 год. URL:
\url{https://rosstat.gov.ru/folder/313/document/279001}.

\bibitem{chi2025}
Chi T.\,C., Fan T.-H., Ghigliazza R.\,M., Giannone D., Wang~Z.
Macroeconomic Forecasting and Machine Learning.
\textit{arXiv:2510.11008}, 2025.

\bibitem{cerqua2025}
Cerqua A., Letta M., Pinto G.
On the (Mis)Use of Machine Learning With Panel Data.
\textit{Oxford Bulletin of Economics and Statistics}, 2025.

\bibitem{badykova2026}
Бадыкова И.\,Р. Валовой региональный продукт: анализ детерминант
для России с использованием моделей машинного обучения.
\textit{Экономика региона}, 2026. Т.~22. №~1. С.~97--108.
DOI: 10.17059/ekon.reg.2026-1-8.

\bibitem{lashina2023}
Lashina M., Grishunin S.
Comparison of Forecasting Power of Statistical Models for GDP Growth
Under Conditions of Permanent Crises.
\textit{Procedia Computer Science}, 2023. Vol.~221. P.~442--449.
DOI: 10.1016/j.procs.2023.07.059.

\bibitem{sigrist2022}
Sigrist F. Gaussian Process Boosting.
\textit{Journal of Machine Learning Research}, 2022.
Vol.~23. No.~232. P.~1--46.

\bibitem{smirnov2025}
Смирнов С.\,В. Методы машинного обучения в макроэкономическом
прогнозировании: предварительные итоги.
\textit{Вопросы экономики}, 2025. №~10. С.~131--154.
DOI: 10.32609/0042-8736-2025-10-131-154.

\bibitem{tsukarev2025}
Цукарев Т., Погосян К., Лемба К., Гришин Д.
Наукастинг ВВП: от традиционных эконометрических моделей к машинному
обучению. \textit{Рабочий документ ЕФСР}, РД/25/1, 2025.
URL: \url{https://efsd.org/research/working-papers/}.

\bibitem{makeeva2022}
Макеева Н.\,М., Станкевич И.\,П.
Наукастинг элементов использования ВВП России.
\textit{Экономический журнал ВШЭ}, 2022. Т.~26. №~4. С.~598--622.
DOI: 10.17323/1813-8691-2022-26-4-598-622.

\bibitem{malkina2025}
Малкина М.\,Ю., Сочков А.\,Л., Капустина Ю.\,И.
Прогнозирование ВРП региона с применением методов искусственного
интеллекта.
\textit{Journal of New Economy}, 2025. Т.~26. №~3. С.~67--85.
DOI: 10.29141/2658-5081-2025-26-3-4.

\bibitem{semiturkin2022}
Semiturkin O., Shevelev A.
Forecasting Regional Inflation Rates Using Machine Learning Methods:
The Case of Siberia Macroregion.
\textit{Bank of Russia Working Paper Series}, No.~91, 2022.

\bibitem{chernyshova2022}
Chernyshova G.\,Yu., Rasskazkin N.\,D.
Software Implementation of Ensemble Models for the Analysis of
Regional Socio-Economic Development Indicators.
\textit{Izvestiya of Saratov University. Mathematics. Mechanics.
Informatics}, 2022. Vol.~22. Iss.~1. P.~130--137.
DOI: 10.18500/1816-9791-2022-22-1-130-137.

\bibitem{buckmann2025}
Buckmann M., Potjagailo G.
Infusing Economically Motivated Structure into Machine Learning Methods.
\textit{Bank of England Staff Working Paper}, 2025.

\bibitem{attolico2025}
Attolico L.
Explainable Machine Learning for Macroeconomic and Financial
Nowcasting. \textit{arXiv:2512.00399}, 2025.

\bibitem{lundberg2017}
Lundberg S.\,M., Lee S.-I.
A Unified Approach to Interpreting Model Predictions.
\textit{Advances in Neural Information Processing Systems}, 2017.

\bibitem{breiman2001}
Breiman L. Random Forests.
\textit{Machine Learning}, 2001. Vol.~45. P.~5--32.

\bibitem{ke2017}
Ke G., Meng Q., Finley T., et al.
LightGBM: A Highly Efficient Gradient Boosting Decision Tree.
\textit{Advances in Neural Information Processing Systems}, 2017.

\bibitem{dorogush2018}
Dorogush A.\,V., Ershov V., Gulin A.
CatBoost: Gradient Boosting with Categorical Features Support.
\textit{arXiv:1810.11363}, 2018.

\end{thebibliography}

\end{document}
